% Generated draft: requires mathematical and source review.
\documentclass[11pt]{article}
\usepackage[margin=25mm]{geometry}
\usepackage{fontspec}
\setmainfont{Latin Modern Roman}
\usepackage{amsmath,amssymb,mathtools,unicode-math}
\setmathfont{Latin Modern Math}
\usepackage{enumitem,needspace,xurl,hyperref}
\setlength{\parindent}{0pt}
\setlength{\parskip}{5pt}
\setlength{\emergencystretch}{3em}
\newcommand{\N}{\mathbb{N}}
\newcommand{\Z}{\mathbb{Z}}
\newcommand{\Q}{\mathbb{Q}}
\newcommand{\R}{\mathbb{R}}
\newcommand{\C}{\mathbb{C}}
\newcommand{\F}{\mathbb{F}}
\newcommand{\A}{\mathbb{A}}
\newcommand{\PP}{\mathbb{P}}
\newcommand{\E}{\mathbb{E}}
\newcommand{\eps}{\varepsilon}
\newcommand{\dd}{\,\mathrm{d}}
\begin{document}
\section*{UnsolvedMath ID: 30001870}
\textbf{Source problem code:} catalogue entry
\subsection{Definitions and mathematical statement}
Let \(X\) be a projective variety over a field \(K\), and write \(n=\dim X\). Under the interpretation that an integral divisor means a Cartier divisor with integer coefficients, define
\[
 h^0(X,mD):=\dim_K H^0\bigl(X,\mathcal O_X(mD)\bigr),
 \qquad
 \operatorname{vol}_X(D):=
 \limsup_{\substack{m\to\infty\\m\in\Z_{>0}}}
 \frac{n!\,h^0(X,mD)}{m^n}.
\]
Here \(\mathcal O_X(mD)\) is the invertible sheaf associated with the Cartier divisor \(mD\).

The supplied realizability question is: does every \(r\in\R_{\geq 0}\) admit a projective variety \(X\) and an integral divisor \(D\) on \(X\) such that
\[
 \operatorname{vol}_X(D)=r?
\]
Both \(X\) and \(D\) may depend on \(r\); the source fixes no dimension and imposes no explicit smoothness, normality, ampleness, or bigness hypothesis.

The supplied record does not specify a ground field, whether that field is fixed throughout the question, or whether ``integral divisor'' means an integral Cartier divisor or an integral Weil divisor. Consequently, the displayed definition gives the Cartier interpretation, while the admissible class of varieties and divisors remains unresolved. A Weil-divisor interpretation would require additional hypotheses or conventions for the spaces of sections. These choices must be checked against the source before this becomes a fully specified mathematical statement.
\subsection{Short English statement}
The question asks whether every nonnegative real number can be realized as the asymptotic volume of an integral divisor on a projective variety, allowing the variety and its dimension to vary. The ground field and the precise meaning of integral divisor are unspecified in the supplied record.
\subsection{Sources}
\begin{itemize}
\item Supplied local UnsolvedMath registry record: permanent ID \(30001870\), title ``Realizability of Divisor Volumes,'' source problem code \texttt{catalogue entry}, in the Oberwolfach Reports Open Problems collection. The record gives proposed year \(2011\) and supplies no named proposer. It labels the statement ``exact'' and reports a comparison with the official report; that verification is supplied provenance, not an independent verification here.
\item Supplied source citation, also listed as ``original OWR report'': \emph{Mini-Workshop: New Developments in Newton-Okounkov Bodies (2012)}, Oberwolfach Reports, pp.~2327--2363, DOI \texttt{10.4171/owr/2011/41}. \url{https://doi.org/10.4171/owr/2011/41}. The background describes these reports as records of mathematical workshops and their open-problem discussions.
\item Supplied literature source: \emph{Convex bodies associated to linear series}. \url{https://arxiv.org/abs/0805.3885}.
\item Supplied literature source: \emph{asymptotic invariants of line bundles}. \url{https://arxiv.org/abs/math/0505054}.
\item Literature assessment dated exactly \texttt{2026-08-21}: the supplied curation labels the problem ``partially solved.'' It states that Riemann--Roch and asymptotic volume theory provide substantial asymptotic information, but that no uniform formula for \(h^0(X,kD)\) on arbitrary projective varieties with Cartier divisors is known. This is supplied dated curation, not independently verified; the record explicitly requests independent verification of time-sensitive status and cited claims. Its relevance to the precise volume-realizability question requires review.
\end{itemize}
\end{document}
